% Options for packages loaded elsewhere
\PassOptionsToPackage{unicode}{hyperref}
\PassOptionsToPackage{hyphens}{url}
\PassOptionsToPackage{dvipsnames,svgnames,x11names}{xcolor}
\PassOptionsToPackage{space}{xeCJK}
\documentclass[
]{article}
\usepackage{xcolor}
\usepackage[margin=2.5cm]{geometry}
\usepackage{amsmath,amssymb}
\setcounter{secnumdepth}{-\maxdimen} % remove section numbering
\usepackage{iftex}
\ifPDFTeX
  \usepackage[T1]{fontenc}
  \usepackage[utf8]{inputenc}
  \usepackage{textcomp} % provide euro and other symbols
\else % if luatex or xetex
  \usepackage{unicode-math} % this also loads fontspec
  \defaultfontfeatures{Scale=MatchLowercase}
  \defaultfontfeatures[\rmfamily]{Ligatures=TeX,Scale=1}
\fi
\usepackage{lmodern}
\ifPDFTeX\else
  % xetex/luatex font selection
    \setmainfont[]{DejaVu Sans}
    \setmonofont[]{DejaVu Sans Mono}
  \ifXeTeX
    \usepackage{xeCJK}
    \setCJKmainfont[]{Noto Sans CJK SC}
          \fi
  \ifLuaTeX
    \usepackage[]{luatexja-fontspec}
    \setmainjfont[]{Noto Sans CJK SC}
  \fi
\fi
% Use upquote if available, for straight quotes in verbatim environments
\IfFileExists{upquote.sty}{\usepackage{upquote}}{}
\IfFileExists{microtype.sty}{% use microtype if available
  \usepackage[]{microtype}
  \UseMicrotypeSet[protrusion]{basicmath} % disable protrusion for tt fonts
}{}
\makeatletter
\@ifundefined{KOMAClassName}{% if non-KOMA class
  \IfFileExists{parskip.sty}{%
    \usepackage{parskip}
  }{% else
    \setlength{\parindent}{0pt}
    \setlength{\parskip}{6pt plus 2pt minus 1pt}}
}{% if KOMA class
  \KOMAoptions{parskip=half}}
\makeatother
\usepackage{longtable,booktabs,array}
\usepackage{calc} % for calculating minipage widths
% Correct order of tables after \paragraph or \subparagraph
\usepackage{etoolbox}
\makeatletter
\patchcmd\longtable{\par}{\if@noskipsec\mbox{}\fi\par}{}{}
\makeatother
% Allow footnotes in longtable head/foot
\IfFileExists{footnotehyper.sty}{\usepackage{footnotehyper}}{\usepackage{footnote}}
\makesavenoteenv{longtable}
\setlength{\emergencystretch}{3em} % prevent overfull lines
\providecommand{\tightlist}{%
  \setlength{\itemsep}{0pt}\setlength{\parskip}{0pt}}
\usepackage{bookmark}
\IfFileExists{xurl.sty}{\usepackage{xurl}}{} % add URL line breaks if available
\urlstyle{same}
\hypersetup{
  colorlinks=true,
  linkcolor={Maroon},
  filecolor={Maroon},
  citecolor={Blue},
  urlcolor={Blue},
  pdfcreator={LaTeX via pandoc}}

\author{}
\date{}

\begin{document}

\section{统一框架中的三通道等价：构造获得直觉}\label{ux7edfux4e00ux6846ux67b6ux4e2dux7684ux4e09ux901aux9053ux7b49ux4ef7ux6784ux9020ux83b7ux5f97ux76f4ux89c9}

\section{Three-Channel Equivalence in the Unified Framework of
Representation, Logic, and Intuition --- Construction Earns
Intuition}\label{three-channel-equivalence-in-the-unified-framework-of-representation-logic-and-intuition-construction-earns-intuition}

\begin{quote}
\textbf{中文版 v0.1} \textbar{} 2026-08-11 (翻译自英文版) 目标: TMLR /
COLM / Cognition (理论 + 实证) \textbf{框架定位}:
属于《表示、逻辑、直觉的统一框架》(纲领:
\texttt{Unified\_Framework\_Representation\_Logic\_Intuition.md})。本文确立
\textbf{构造→直觉} 的认知环节:
构造、直觉、形式重写是同一系统的三个等价执行通道;
直觉是编译后的构造。姊妹篇: P1 (形式→构造, 表示 + 泛化理论), P2
(构造→直觉方法, 编译)。 证据: exp50 (三通道一致), exp41 (符号置换),
exp80 (跳步), exp53 (收敛)。单一官方 run\_exp 判定。
\end{quote}

\begin{center}\rule{0.5\linewidth}{0.5pt}\end{center}

\textbf{Repository:
\href{https://github.com/YuchenWang-ai/Three-Channel-Equivalence-Construction-Intuition-and-Formal-Rewriting-in-One-Neural-System}{YuchenWang-ai/Three-Channel-Equivalence-Construction-Intuition-and-Formal-Rewriting-in-One-Neural-System}}

\subsection{摘要}\label{ux6458ux8981}

数学基础中长达一个世纪的分裂 --- 直觉主义
(数学对象存在当且仅当它能被构造)、结构主义
(对象即其在结构中的位置)、形式主义 (数学是一场符号游戏) ---
在单一神经系统中镜像为同一数学对象的三个执行通道。我们证明: 在 token
原生表示中, \emph{构造} 通道 (定义驱动展开)、\emph{直觉} 通道
(编译后的权重执行)、\emph{形式} 通道 (定义驱动重写) 两两语义等价: 在
2000 位进位链上三者产生完全相同的结果 (1.000)。我们论证这是计算三位一体
(Computational Trinitarianism, Curry--Howard: 程序 = 证明 = 类型)
的神经实现, 关键区别在于等价发生在\textbf{同一系统的执行通道之间},
而非不同逻辑系统之间。我们发展认识论: 直觉是推理快路径 --- 难以获得
(配置相关的轮数), 但易于通过跳过构造步骤而泛化; 符号置换不变性 (0.987)
与零衰减跳步 (2000 位、60 进制、100\% 匿名化)
确立直觉通道承载的是结构而非符号语义。立场: 反柏拉图主义 ---
直觉是编译后的构造, 而非先天能力。在统一框架内, 本文是
\textbf{构造→直觉} 的认知环节: 正确的构造产生直觉 (编译的构造),
正确的泛化编译为闪念直觉 (快速推理路径)。

\subsection{1 引言}\label{ux5f15ux8a00}

\textbf{分裂。} 在数学基础中,
三个答案竞争''数学对象是什么''这一问题。布劳威尔的直觉主义:
对象存在当且仅当它能被\emph{构造}。布尔巴基的结构主义:
对象即其在结构中的\emph{位置}。希尔伯特的形式主义:
数学是一场\emph{符号游戏};
对象即其在重写系统中的角色。这一争论通常被呈现为哲学问题;
在本文中我们观察到它同时也是\emph{架构问题} ---
三种立场对应计算系统实现同一对象的三种方式。

\textbf{统一。} 一个基于 token 原生表示 (P1)
的单一神经系统中存在全部三个通道: - \emph{构造} (直觉主义):
逐步展开定义链 --- 由构造而存在; - \emph{直觉} (结构主义):
编译后的权重直接定位对象 --- 由在学习结构中的位置而存在; - \emph{形式}
(形式主义): 定义驱动的重写引擎 --- 由符号操作规则而存在。

我们在实验中证明三个通道在相同输入上两两语义等价 (exp50: 构造 =
形式完全一致; 直觉在 2000 位进位链上 1.000)。这是计算三位一体的神经实现
--- Curry--Howard 的 程序 = 证明 = 类型 --- 但有一个决定性区别:
等价发生在\emph{同一系统的执行通道之内}, 而非不同逻辑系统之间。

\textbf{为什么重要。} 如果构造、直觉、形式重写对同一对象给出相同答案,
那么这些哲学立场就不再是需要争论的备选项,
而是\emph{单一实现的不同侧面}。``对象是什么''的分歧变成''透过哪个通道看''的分歧
--- 而所有通道给出相同的回答。工程内容在于编译关系: 构造可以被编译为直觉
(P2), 而两者都与重写一致。

\textbf{贡献。}

\begin{itemize}
\tightlist
\item
  \textbf{C1 (认识论) --- 直觉是编译后的构造。}
  直觉通道不是第三种神秘能力; 它是被编译进权重的构造过程。反柏拉图主义:
  无先天结构, 无先验对象。
\item
  \textbf{C2 (实证) --- 三通道等价。} 构造 = 形式完全一致,
  直觉在相同输入上 1.000, 包括 2000 位进位链 (exp50)。
\item
  \textbf{C3 (实证) --- 直觉承载结构而非符号。} 符号置换不变性 (0.987,
  exp41) 与零衰减跳步 (exp80) 表明直觉通道实现的是对象的结构,
  而非其能指。
\item
  \textbf{C4 (认识论) --- 获得与跳步解耦。} 直觉难以获得 (配置相关轮数;
  简单 3 轮 / 完整 8 轮, exp53) 但易于通过跳过构造步骤而泛化 ---
  两者并不矛盾。
\end{itemize}

\subsection{2 三种立场,
三个通道}\label{ux4e09ux79cdux7acbux573a-ux4e09ux4e2aux901aux9053}

\begin{longtable}[]{@{}lllll@{}}
\toprule\noalign{}
立场 & 对象是 & 访问方式 & 我们的通道 & 执行 \\
\midrule\noalign{}
\endhead
\bottomrule\noalign{}
\endlastfoot
直觉主义 (布劳威尔) & 一个构造 & 构造它 & 构造 & 展开定义链, 逐步 \\
结构主义 (布尔巴基) & 一个位置 & 定位它 & 直觉 & 编译权重直接放置 \\
形式主义 (希尔伯特) & 一场符号游戏 & 玩它 & 形式 & 定义驱动重写 \\
\end{longtable}

三个通道共享单一定义源: token 原生表示 (P1
§4)。构造与形式读取相同的定义;
直觉是同一定义的编译映像。因此等价不是三个独立系统的巧合,
而是单一系统透过三个通道观察到的性质。

\textbf{计算三位一体。} Curry--Howard
将程序、证明、类型视为同一事物的三种呈现。我们的主张类似:
构造、直觉、形式重写是同一计算在三种呈现下的形态 ---
\emph{在单一系统之内}。与 Curry--Howard 的区别在于等价的所在: 不是''程序
= 证明 = 类型''作为不同实体, 而是''同一对象的三个执行路径一致''。

\subsection{3
直觉作为编译的构造}\label{ux76f4ux89c9ux4f5cux4e3aux7f16ux8bd1ux7684ux6784ux9020}

\textbf{主张。} 直觉通道是被编译进权重的构造过程。训练 (P1 §5)
用判定结果监督模型, 同时\emph{排除答案轨迹}: 模型无法复制展开过程;
它必须从定义中归纳计算并编译之。编译产物即直觉 --- 推理快路径。

\textbf{难以获得。} 编译不是免费的。获得是配置相关的: 简单套件 3 轮,
完整逻辑+算术套件 8 轮 (exp53)。编译完成前, 直觉通道不泛化 (exp53: 第
1--3 轮 0.1--24\%)。这是解耦的\emph{获得}侧。

\textbf{易于泛化。} 一旦编译完成, 直觉通道通过\emph{跳过构造步骤}泛化:
它从输入直接跳到输出, 越过中间展开。外推到 2000 位、60 进制、100\%
数字匿名化零衰减 (exp80)。这是\emph{泛化}侧。

\textbf{解耦 (C4)。} 两个事实 --- 难以获得、易于跳步 ---
并不矛盾。获得需要编译结构 (训练); 泛化使用编译后的结构 (推理)。P1
的理论解释两者: 定义完备 + 关系充分决定编译是否成功;
跳步正是编译结构在新输入上的行为。

\subsection{4 实验}\label{ux5b9eux9a8c}

全部经官方 run\_exp 判定路径 (完整序列重建)。模型: exp02\_supervised\_s2
(2 层, dim-64, \textless1 MB, 完整套件判定 1.000)。

\subsubsection{4.1 三通道等价
(C2)}\label{ux4e09ux901aux9053ux7b49ux4ef7-c2}

\begin{longtable}[]{@{}llll@{}}
\toprule\noalign{}
输入 & 构造 & 形式 & 直觉 \\
\midrule\noalign{}
\endhead
\bottomrule\noalign{}
\endlastfoot
12+34 & ✓ & ✓ & ✓ \\
123456789+987654321 & ✓ & ✓ & ✓ \\
999\ldots9 (20 位)+1 & ✓ & ✓ & ✓ \\
\textbf{999\ldots9 (2000 位)+1} & ✓ & ✓ & \textbf{1.000} \\
\end{longtable}

构造 = 形式在完整输出序列上完全一致 (2/9/20/2000 位)。直觉在相同输入上
1.000。三个通道实现同一计算。

\subsubsection{4.2 直觉承载结构而非符号
(C3)}\label{ux76f4ux89c9ux627fux8f7dux7ed3ux6784ux800cux975eux7b26ux53f7-c3}

\begin{itemize}
\tightlist
\item
  \textbf{符号置换 (exp41)}: 训练中 26080 次数字重命名, 原始符号上测试
  --- 不变性 0.996→0.987。直觉通道实现结构关系 (后继/迭代/进位),
  而非数字能指。
\item
  \textbf{跳步零衰减 (exp80)}: 2000 位 / 60 进制 / 100\% 匿名化 / 联合,
  全部 1.000。编译结构外推无衰减。
\item
  \textbf{机器形式 (exp90)}: 20 位加法将 82 个构造 token 编译为 5 个直觉
  token --- 跳步在 token 流中可见。诚实警示: 原子 token 在模型词表外
  (pad id 0); 测量的是序列长度效应, 为下界。
\end{itemize}

\subsubsection{4.3 获得--泛化解耦
(C4)}\label{ux83b7ux5f97ux6cdbux5316ux89e3ux8026-c4}

\begin{longtable}[]{@{}llllll@{}}
\toprule\noalign{}
轮数 & 1 & 2 & 3 & 5 & 8 \\
\midrule\noalign{}
\endhead
\bottomrule\noalign{}
\endlastfoot
判定准确率 (完整套件) & 0.001 & 0.214 & 0.244 & 0.781 & 0.996 \\
\end{longtable}

获得是单调但配置相关的 (完整套件 8 轮)。同一收敛模型随后零衰减外推
(§4.2)。``难以获得''与''易于跳步''共存; 它们分别描述获得 (训练) 与泛化
(推理)。

\subsection{5 讨论}\label{ux8ba8ux8bba}

\textbf{等价不意味着什么。} 我们不声称三种哲学立场在认识论上相同;
我们声称它们作为\emph{同一系统的执行通道}等价。历史上的争论关乎数学真理的本性;
我们的主张关乎实现数学对象的计算系统的架构。哲学立场是关于数学对象是什么的模型;
我们展示一个系统能一致地实现同一对象的全部三个模型。

\textbf{反柏拉图主义。} 直觉是编译的构造,
而非先天能力。如果数学结构是先天 (柏拉图主义/天赋主义),
那么定义贫乏的概念仍应泛化; 事实并非如此 (P1 §8:
贫乏词表需要高一个数量级的训练; 真正欠定义的 token
永不泛化)。直觉通道的能力由编译赢得, 而非先天赋予。

\textbf{位置内容 vs 构造内容。} 结构主义说对象即其位置;
直觉通道在学习结构中定位对象。直觉主义说对象即其构造;
构造通道展开它。通道的等价表明两种描述一致 ---
对于定义完备的对象。这是结构主义--直觉主义旧日和解的机器级形式。

\subsection{6 相关工作}\label{ux76f8ux5173ux5de5ux4f5c}

\begin{itemize}
\tightlist
\item
  \textbf{计算三位一体} (Curry--Howard): 程序 = 证明 =
  类型。我们将三合一识别扩展到单一神经系统的执行通道。
\item
  \textbf{构造数学 / 类型论} (Martin-Löf 1984; Bishop): 存在 =
  可构造性。我们的构造通道是其机器形式。
\item
  \textbf{结构主义 / 布尔巴基; Gärdenfors 概念空间}: 对象 =
  结构/几何中的位置。我们的直觉通道是其神经实现。
\item
  \textbf{形式主义 / 希尔伯特; 重写系统}: 对象 =
  符号操作。我们的形式通道是定义驱动的引擎。
\item
  \textbf{双过程认知} (Kahneman): 系统 1/2。我们的直觉/推理分离即系统
  1/2 之分, 并给出系统 1 如何被构建的编译解释。
\item
  \textbf{数学认知中的双系统} (如数感 / 刻意算术):
  我们的三通道以可验证的等价细化了双系统图景。
\end{itemize}

\subsection{6.5 数学领域验证:
黎曼方向案例}\label{ux6570ux5b66ux9886ux57dfux9a8cux8bc1-ux9eceux66fcux65b9ux5411ux6848ux4f8b}

三通道的论述不限于算术。在姊妹形式化 (Lean 4 / mathlib, claims
C011--C025, 全部 PROVED, novelty KNOWN, 无 sorry) 中, 直觉链''复轴是 −1
的高维投影; 素数落在平移的整数点上; 1/2 是反转--平移对称中心;
复轴被卷曲''被贯穿全部三个通道:

\begin{itemize}
\tightlist
\item
  \textbf{构造} (定义驱动): ComplexAxis 代数 (a + bJ, J² = −1), 投影
  proj⟨a,b⟩ = a, 提升, 基点漂移;
\item
  \textbf{形式} (重写): Lean 4 证明投影不保持、素数圆结构 (旋转 ×
  共轭下单轨道的 8 个格点, 经高斯整数 UFD)、欧拉乘积收敛 (Re(s)
  \textgreater{} 1)、临界线几何;
\item
  \textbf{直觉} (编译结构): 同一对象经直觉链直接到达, 跳过教科书推导。
\end{itemize}

\textbf{投影恢复定理} (\texttt{proj\_not\_recoverable} /
\texttt{proj\_recoverable\_symmetry}, Lean 证明):
\emph{丢失的结构不可恢复; 对称方向可恢复。} 投影下, 虚轴方向 (i 与 −i)
丢失且不可恢复, 而实轴符号对称 (r 与 −r) 被保持。这是 EXP-61h (§4)
的数学对应: 直觉通道位置不敏感 (它丢弃位置细节, 如投影丢弃 J 方向)
却保持语义对称 --- 与''打乱的字母不妨碍理解;
对打乱内容刻意逻辑反而失去意义'' (G5b, P1) 一致。投影是快路径的机制:
不可逆地丢弃无关结构正是直接路径廉价的原因。

\textbf{直觉引导形式化的 token 经济} (方法论数据): 完整 C011--C025
形式化消耗 ≈700k token, 其中 99.2\% 为上下文缓存重发; 净新内容
\textless{} 1\%。直觉链直接命中已知结构
(heap/torsor、两平方和、圆反转、函数方程对称), 避开教科书推导 --- 这是
G5 在 token 层面的效率数据, 独立于 §4 的算术实验。

\subsection{7 局限}\label{ux5c40ux9650}

\begin{enumerate}
\def\labelenumi{\arabic{enumi}.}
\tightlist
\item
  等价仅在受支持演算 (形式化定义) 上演示, 而非全部数学。
\item
  直觉通道是统计性的, 非保证; 等价是验证 (留出集) 的并回退到构造,
  而非证明。
\item
  原子化结果 (exp90) 是序列长度下界, 而非习得原子语义的主张。
\end{enumerate}

\subsection{8 结论}\label{ux7ed3ux8bba}

基础之争的三种立场 --- 构造、结构、符号 ---
被实现为单一神经系统的三个等价执行通道。构造 = 形式完全一致; 直觉 1.000
一致; 直觉通道承载结构而非符号;
它难以获得但易于跳步泛化。直觉是编译的构造。数学基础中一个世纪的分裂,
至少对一个形式化演算而言,
是''一个系统能以多少种方式实现同一对象''的架构事实 --- 而它们全部一致。

\subsection{附录}\label{ux9644ux5f55}

\begin{itemize}
\tightlist
\item
  A. 完整实验表: docs/paper\_data/results.csv。
\item
  B. 与 P1 命题 (G0, G5, R) 及 P2 编译的关系。
\item
  C. 形式验证方向 (Lean): 构造 ≡ 直觉 ≡ 形式, 对受支持演算。
\end{itemize}

\begin{center}\rule{0.5\linewidth}{0.5pt}\end{center}

\subsection{统一结论 (2026-08-11
增补)}\label{ux7edfux4e00ux7ed3ux8bba-2026-08-11-ux589eux8865}

当我们在统一的形式化注意力语法结构下,
最终不可避免地、必然地通过神经网络清晰、直观、可解释地观察到直觉与构造的边界时,
我们就可以有充足的信心: 用统一的注意力形式语言指导逻辑构造, 用精确的合成
CoT 逻辑构造获得 \emph{闪念直觉},
用蒸馏的直觉上下文穷举注意力形式。形式、构造、直觉在统一框架下相互适配,
组成再不分彼此的教学、训练、推理的反馈迭代管线 --- 这也许是我们通往 ASI
的众多路径中, 相对较快的一条捷径。

\subsection{统一框架内的交叉引用}\label{ux7edfux4e00ux6846ux67b6ux5185ux7684ux4ea4ux53c9ux5f15ux7528}

本文是《表示、逻辑、直觉的统一框架》(纲领:
\texttt{Unified\_Framework\_Representation\_Logic\_Intuition.md})
的一条腿。三条腿相互适配:

\begin{longtable}[]{@{}lll@{}}
\toprule\noalign{}
腿 & 论文 & 主张 \\
\midrule\noalign{}
\endhead
\bottomrule\noalign{}
\endlastfoot
形式 → 构造 & P1 (本文) & 正确语法 → 极快训练; 正确构造 → 极强泛化 \\
构造 → 直觉 & P0 (三通道等价) & 构造 = 形式 = 直觉; 直觉是编译的构造 \\
构造 → 直觉 (方法) & P2 (神经宏编译) & 把慢构造路径编译为快推理路径 \\
\end{longtable}

循环闭合: 形式指导构造 (P1), 构造获得直觉 (P0/P2), 直觉搜索形式
(快路径外推并再表达于统一形式)。统一结论 --- 形式指导逻辑构造,
构造获得闪念直觉, 直觉穷举搜索形式 ---
在纲领与每篇论文的统一结论中陈述。

\end{document}
